\caption{Formal three-seed sequential diagnostics on ImageNet-A. From left to right: effective Q/V weight drift from each method's own Task0 anchor; signed Task0 test cross-entropy change relative to that anchor; Task0 test accuracy; and normalized directional path sensitivity (``Average Fisher'') $\overline S_{\mathrm{path},t}=\int_0^1\Delta_t^\top F_0^{\mathrm{emp}}(W_0+s\Delta_t)\Delta_t\,ds/\|\Delta_t\|_2^2$, with $\Delta_t=W_t-W_0$. The horizontal coordinate is the number of learned tasks $T=t+1$. The first three panels contain all ten measured endpoints $T=1,\ldots,10$. Loss and accuracy use all 307 old-task test examples and the fixed 20-class Task0 head. The last panel contains independently measured anchor-to-endpoint straight chords for every $T=2,\ldots,10$, with no estimate at $T=1$ because the denominator is zero. Each chord's empirical Fisher uses all 1,259 old-task training examples, true-label per-example gradients in the 24 effective Q/V matrices, and no EWC coefficient, trace normalization, floor, or regularization. Adaptive nested Simpson quadrature used 9, 17 measured nodes per chord; all final stability and numerical-bound checks passed. Integration nodes are positions along a chord, not task stages; these chords are not the optimizer's chronological path. Curves show arithmetic means over training seeds 0, 1, and 2, and shaded bands show sample standard deviations ($n=3$, denominator $n-1$). Sensitivity uses each seed's own normalized integral before aggregation and is not an endpoint Fisher trace. The original 18 paths at $t=1,5,9$ are reused; 36 new paths fill $t=2,3,4,6,7,8$.}
